\documentclass{article}
\usepackage[a4paper, margin=1in]{geometry}
\usepackage{amsmath, amssymb}
\usepackage{enumitem}
\usepackage{hyperref}

\title{}
\date{}

\begin{document}

\maketitle

\section*{Semi-total Domination in Unit Disk Graphs}


\textbf{Authors:} Sasmita Rout, Gautam Kumar Das

\textbf{Published in:} CALDAM 2024 (LNCS 14508)

\textbf{DOI:} \href{https://doi.org/10.1007/978-3-031-52213-0_9}{10.1007/978-3-031-52213-0\_9}

\vspace{1em}

\section*{1. Problem Studied}
\begin{itemize}[leftmargin=1.5em]
    \item The paper studies the \textbf{Semi-total Dominating Set (STDS)} problem in Unit Disk Graphs (UDGs).
    \item A set \( D_{t2} \subseteq V \) is a STDS if:
    \begin{enumerate}
        \item \( D_{t2} \) is a dominating set.
        \item For every \( u \in D_{t2} \), there exists \( v \in D_{t2} \) such that \( \text{dist}(u,v) \leq 2 \).
    \end{enumerate}
\end{itemize}

\section*{2. Motivation and Applications}
\begin{itemize}[leftmargin=1.5em]
    \item STDS is a natural variant lying between classical dominating sets and total dominating sets.
    \item Relevant in wireless networks and sensor placements where both coverage and minimal connectivity are required.
\end{itemize}

\section*{3. Main Contributions}
\begin{itemize}[leftmargin=1.5em]
    \item Proved that the STDS problem is \textbf{NP-complete} in UDGs via a reduction from Vertex Cover on planar graphs.
    \item Proposed a \textbf{6-factor approximation algorithm} for STDS in UDGs.
    \item Introduced a \textbf{2 + ln(D+1)} approximation for general graphs, improving over the previous bound of \textbf{2 + 3ln(D+1)}.
    \item Showed that a PTAS for STDS in UDGs exists with an approximation factor of \( 2 + \varepsilon \).
\end{itemize}

\section*{4. Techniques Used}
\begin{itemize}[leftmargin=1.5em]
    \item \textbf{NP-Completeness:} Used grid embedding and geometric gadget constructions to simulate vertex cover behavior.
    \item \textbf{Approximation Algorithm:}
    \begin{enumerate}
        \item Compute a maximal independent set \( D \).
        \item For each vertex in \( D \) lacking a neighbor within distance 2 in \( D \), add a neighbor from \( N(u) \) to set \( T \).
        \item Return \( D \cup T \) as the semi-total dominating set.
    \end{enumerate}
\end{itemize}

\section*{5. Results Summary}

\begin{tabular}{|l|l|l|l|}
\hline
\textbf{Graph Type} & \textbf{Approach} & \textbf{Approx. Factor} & \textbf{Time Complexity} \\
\hline
UDG & Greedy + Local Fix & 6 & \( O(nk) \) \\
UDG (PTAS) & Local Search + Packing & \( 2 + \varepsilon \) & \( n^{O(1/\varepsilon \log \varepsilon)} \) \\
General Graph & Set Cover-Based & \( 2 + \ln(D + 1) \) & --- \\
\hline
\end{tabular}

\vspace{1em}

\section*{6. Comparison with Prior Work}
\begin{itemize}[leftmargin=1.5em]
    \item Improves over the \( 2 + 3 \ln(D+1) \) result from Henning and Pandey (2019).
    \item The first work to give a dedicated algorithm and complexity result for STDS in UDGs.
    \item Relates to: \( \gamma(G) \leq \gamma_{t2}(G) \leq \gamma_t(G) \).
\end{itemize}

\section*{7. Limitations / Assumptions}
\begin{itemize}[leftmargin=1.5em]
    \item Assumes no isolated vertices.
    \item Requires geometric embedding (disk centers must be known).
    \item Maximal independent set construction may not be efficient for very large or distributed settings.
\end{itemize}

\section*{8. Observations and Future Ideas}
\begin{itemize}[leftmargin=1.5em]
    \item Potential to explore distributed or dynamic versions of the 6-approximation algorithm.
    \item LP relaxation and rounding might yield tighter bounds.
    \item STDS in other geometric graphs or higher dimensions could be investigated.
\end{itemize}

\section*{(Independent) Roman Domination Parameterized by Distance to Cluster}

\textbf{Authors:} Pradeesha Ashok, Gautam K. Das, Arti Pandey, Kaustav Paul, Subhabrata Paul

\textbf{Published in:} COCOA 2024 (LNCS 15435)

\textbf{DOI:} \href{https://doi.org/10.1007/978-981-96-4448-3_13}{10.1007/978-981-96-4448-3\_13}

\vspace{1em}

\section*{1. Problem Studied}
\begin{itemize}[leftmargin=1.5em]
    \item The paper studies the Roman Domination and Independent Roman Domination problems.
    \item A Roman Dominating Function (RDF) is a function \( f : V \rightarrow \{0,1,2\} \), such that every vertex \( v \) with \( f(v) = 0 \) has a neighbor \( u \) with \( f(u) = 2 \).
    \item An Independent RDF (IRDF) further requires that the set of vertices assigned 1 or 2 form an independent set.
    \item The study focuses on structural parameterization: specifically, the parameter is the size \( k \) of a minimum set of vertices whose removal results in a cluster graph (a disjoint union of cliques).
\end{itemize}

\section*{2. Motivation and Applications}
\begin{itemize}[leftmargin=1.5em]
    \item Roman Domination arises from defense strategies inspired by the Roman Empire.
    \item It models efficient resource placement (e.g., network monitoring, facility location) under constraints.
    \item Independent RDF ensures that no two deployed resources are adjacent — useful in scenarios requiring non-interference.
\end{itemize}

\section*{3. Main Contributions}
\begin{itemize}[leftmargin=1.5em]
    \item Proved that both the Roman Domination (RD-CVD) and Independent Roman Domination (IRD-CVD) problems are \textbf{FPT} when parameterized by distance to cluster.
    \item Gave algorithms with time complexity \( 4^k n^{\mathcal{O}(1)} \) for both problems.
    \item Proved matching \textbf{lower bounds} under the SETH assumption: the problems cannot be solved in \( 2^{\varepsilon k} n^{\mathcal{O}(1)} \) for any \( \varepsilon < 1 \), unless SETH fails.
    \item Proved that both problems do not admit a \textbf{polynomial kernel} unless $NP \subseteq coNP/poly$.
\end{itemize}

\section*{4. Techniques Used}
\begin{itemize}[leftmargin=1.5em]
    \item Used \textbf{structured branching} on cluster vertex deletion sets.
    \item Reduced the problem to a variant of \textbf{Set Cover with partition constraints}.
    \item Designed a dynamic programming algorithm with state space bounded by \( 2^{|S|} \cdot n^{\mathcal{O}(1)} \), where \( S \) is the cluster deletion set.
    \item Used reductions from the \textbf{d-Hitting Set} problem to prove lower bounds and kernelization hardness.
\end{itemize}

\section*{5. Results Summary}
\begin{tabular}{|l|l|l|l|}
\hline
\textbf{Problem} & \textbf{Parameter} & \textbf{Complexity} & \textbf{Remarks} \\
\hline
RD-CVD & Cluster Deletion Size \( k \) & \( 4^k n^{\mathcal{O}(1)} \) & FPT algorithm \\
IRD-CVD & Cluster Deletion Size \( k \) & \( 4^k n^{\mathcal{O}(1)} \) & FPT algorithm \\
RD-CVD / IRD-CVD & Cluster Deletion Size \( k \) & No \( 2^{\varepsilon k} n^{\mathcal{O}(1)} \) algo & Unless SETH fails \\
RD-CVD / IRD-CVD & Cluster Deletion Size \( k \) & No Poly Kernel & Unless $NP \subseteq coNP/poly$ \\
\hline
\end{tabular}

\section*{6. Comparison with Prior Work}
\begin{itemize}[leftmargin=1.5em]
    \item Extends the line of work on domination problems in parameterized complexity.
    \item Builds upon earlier results on classical Roman Domination and parameterization by treewidth and solution size.
    \item One of the first works to explore Independent Roman Domination with structural parameters.
\end{itemize}

\section*{7. Limitations and Assumptions}
\begin{itemize}[leftmargin=1.5em]
    \item The cluster vertex deletion set is assumed as part of the input.
    \item Results focus on deletion to cluster graphs; extension to other structural parameters like treewidth, neighborhood diversity is open.
    \item Real-world graphs may not always exhibit small cluster distance.
\end{itemize}

\section*{8. Observations and Future Directions}
\begin{itemize}[leftmargin=1.5em]
    \item Strongly bridges parameterized complexity and classical domination theory.
    \item Open directions:
    \begin{itemize}
        \item Improve the time complexity (e.g., reduce base from 4).
        \item Explore FPT status with parameters like modular-width, twin-cover.
        \item Investigate kernelization under alternative parameters.
    \end{itemize}
\end{itemize}

\section*{Total Roman Domination and Total Domination in Unit Disk Graphs}

\textbf{Authors:} Sasmita Rout, Pawan Kumar Mishra, Gautam Kumar Das

\textbf{Published in:} Communications in Combinatorics and Optimization, Vol. 10, No. 4 (2025), pp. 803–823

\textbf{DOI:} \href{https://doi.org/10.22049/cco.2024.28647.1650}{10.22049/cco.2024.28647.1650}

\vspace{1em}

\section*{1. Problem Studied}
\begin{itemize}[leftmargin=1.5em]
    \item The paper focuses on two problems in Unit Disk Graphs (UDGs):
    \begin{enumerate}
        \item The \textbf{Total Roman Dominating Set (TRDS)} problem.
        \item The \textbf{Total Dominating Set (TDS)} problem.
    \end{enumerate}
    \item A \textbf{Total Roman Dominating Function (TRDF)} is a Roman dominating function where the subgraph induced by non-zero valued vertices has no isolated vertex.
    \item The objective is to find a TRDF of minimum weight (sum of vertex values).
\end{itemize}

\section*{2. Motivation}
\begin{itemize}[leftmargin=1.5em]
    \item These problems have direct applications in wireless sensor and ad-hoc networks.
    \item Ensures both domination and resilience (total property), mimicking fault tolerance and mutual communication.
    \item UDGs model proximity-based communication (e.g., sensors with fixed radius).
\end{itemize}

\section*{3. Main Contributions}
\begin{itemize}[leftmargin=1.5em]
    \item Proposed a \textbf{10.5-approximation algorithm} for TRDS in UDGs with time complexity \( \mathcal{O}(|V| \log k) \), where \( k \) is the number of vertices assigned Roman value 2.
    \item Improved approximation factor over previous best-known 12-approximation.
    \item Proposed a new \textbf{7.79-approximation algorithm} for the TDS problem in UDGs, improving over the previous 8-factor approximation.
    \item TDS algorithm runs in \( \mathcal{O}(|V| + |E|) \) time.
\end{itemize}

\section*{4. Techniques Used}
\begin{itemize}[leftmargin=1.5em]
    \item For TRDS:
    \begin{itemize}
        \item Construct maximal independent set \( V_2 \) (assigned value 2).
        \item Add neighboring vertices \( V_1 \) (assigned value 1) to ensure no isolated vertices in \( V_1 \cup V_2 \).
    \end{itemize}
    \item For TDS:
    \begin{itemize}
        \item BFS to layer graph into sets \( S_i \), construct independent sets \( I_i \), and parent sets \( T_i \) for total property.
    \end{itemize}
\end{itemize}

\section*{5. Results Summary}
\begin{tabular}{|l|l|l|l|}
\hline
\textbf{Problem} & \textbf{Approx. Factor} & \textbf{Time Complexity} & \textbf{Improvement} \\
\hline
TRDS in UDGs & 10.5 & \( \mathcal{O}(|V| \log k) \) & Improves previous 12-factor \\
TDS in UDGs & 7.79 & \( \mathcal{O}(|V| + |E|) \) & Improves previous 8-factor \\
\hline
\end{tabular}

\vspace{1em}

\section*{6. Comparison with Prior Work}
\begin{itemize}[leftmargin=1.5em]
    \item Builds on prior work by Shang et al. (RDS in UDGs) and Jena \& Das (TDS in UDGs).
    \item Avoids distributed algorithm overhead from previous CRDS approximations.
    \item Provides improved constants for both TRDS and TDS.
\end{itemize}

\section*{7. Assumptions and Limitations}
\begin{itemize}[leftmargin=1.5em]
    \item Requires UDG embedding (disk centers in \( \mathbb{R}^2 \) known).
    \item Performance guarantees are worst-case; empirical performance not evaluated.
    \item TRDS approximation relies on geometric packing bounds (disk overlaps).
\end{itemize}

\section*{8. Observations and Future Work}
\begin{itemize}[leftmargin=1.5em]
    \item Provides structural lemmas and tight bounds on packing independent disks.
    \item Opens avenues for:
    \begin{itemize}
        \item Distributed versions of the algorithms.
        \item Extending to dynamic or weighted UDGs.
        \item Empirical evaluation on real-world networks.
    \end{itemize}
\end{itemize}

\end{document}


\end{document}
